\scanpage{9}
\uline{Questions pour Artin}

1°) $S$ \deleted{strictement} local \deleted{\unsure{noeth|noeth.}} \uuline{hens.} noeth, $S'\to S$ \deleted{\unsure{fpt.}} plat. surjectif, $S'$ noeth, $X'/S'$ \deleted{\unsure{de p.f.|loc. de t.f.}}\inserted{\uline{propre}} avec donnée de descente rel. à $S'/S$. Supposons que la donnée de descente sur $X'_s$ soit effective ($s$ pt. fermé de $S$). Alors la donnée de descente sur $X'$ est effective ?

Cas où $S'=\widehat{S}$ ?

2°) Peut-on faire la descente fpqc des \uline{variétés} \inserted{utiliser \unsure{théor.|loc.} de pts. fixes,}, p.ex. à partir d'un complété d'un anneau local ? [i.e. la propriété pour un faisceau fpqc loc. de prés. finie d'être une variété est-elle loc. \unsure{fpqc|fppc} sur $S$ ?]

% slanted
\margin{Le Pb est la comm. aux $\underset{\leftarrow}{\lim}$ (\unsure{adiques artiniennes})}

3°) $S$ comme dans \uline{1°}, $X$ une \uline{variété} \deleted{sur, \unsure{loc. de pr. f.|loc. de t.f.}}\inserted{propre} sur $S$. Si $X$ est représentable loc. \unsure{fppf|fppqc|fpqc} sur $S$, \inserted{et $X_s$ représentable,} alors $X$ est-il représentable ? Evidemment, si la réponse à 2° \deleted{\unsure{et 3°|est}} était affirmative, les questions 1° et 3° seraient équivalentes...

% slanted
\margin{\unsure{Exemple|Exemples} de $X$ plat p\unsure{rés.} propre sur $S$ local complet, $X_s$ repr. mais $X$ non repr. ?? Exemple de $X$ \unsure{vari.|variétés} \deleted{\unsure{de}} \unsure{pfppf|propres} sur $S$ noeth, à fibres repr., qui ne soit repr. ?}

4°) Descente fpqc des \inserted{(``petits'')} faisceaux étales.

Lié à la question de \unsure{critériser|conditions sur} les foncteurs $(\mathrm{Sch})_{/S}^{\circ}\to(\mathrm{Ens})$ qui proviennent de \deleted{petits} faisceaux étales sur $S$. Une réponse affirmative à 2° (\unsure{où|si|i.e.} on prendrait ``variété'' dans le sens général, sans exiger que le \inserted{morphisme} \unsure{diagonale|diagonal} soit une immersion) en impliquerait une à 4°.

% archive sheet number: 8
